\section{Lecture 24: Transformers}

\subsection{Motivation}

In the previous lecture, we introduced attention as a mechanism for modeling interactions across sequences.

\medskip

\noindent
\textbf{Key question.}  
Can we build an entire architecture based solely on attention?

\medskip

\noindent
\textbf{Answer.}  
Transformers replace recurrence and convolution with attention, enabling scalable and highly expressive models.

\medskip

\noindent
\textbf{Guiding idea.}  
Representations are computed through repeated layers of attention and simple transformations.

\subsection{Transformer Architecture}

A transformer is composed of stacked layers, each containing:

\begin{itemize}
    \item multi-head self-attention
    \item feedforward network
    \item residual connections
    \item normalization
\end{itemize}

\medskip

\noindent
\textbf{Layer structure.}

Given input $x$:

\[
z = \text{SelfAttention}(x),
\quad
x' = x + z
\]

\[
\tilde{x} = \text{LayerNorm}(x')
\]

\[
u = \text{FFN}(\tilde{x}),
\quad
y = \tilde{x} + u
\]

\[
\text{output} = \text{LayerNorm}(y)
\]

\medskip

\noindent
\textbf{Feedforward network.}

\[
\text{FFN}(x) = \sigma(xW_1 + b_1) W_2 + b_2.
\]

\medskip

\noindent
\textbf{Insight.}  
Transformers alternate between global interaction (attention) and local transformation (FFN).

\subsection{Positional Encoding}

Attention is permutation-invariant, so positional information must be added.

\medskip

\noindent
\textbf{Sinusoidal encoding.}

\[
\text{PE}(t, 2k) = \sin\left(\frac{t}{10000^{2k/d}}\right), \quad
\text{PE}(t, 2k+1) = \cos\left(\frac{t}{10000^{2k/d}}\right).
\]

\medskip

\noindent
\textbf{Purpose.}

\begin{itemize}
    \item encode order information
    \item allow model to distinguish positions
\end{itemize}

\medskip

\noindent
\textbf{Insight.}  
Transformers separate content (tokens) from position (encoding).

\subsection{Residual Connections and Normalization}

Transformers rely heavily on:

\begin{itemize}
    \item residual connections
    \item layer normalization
\end{itemize}

\medskip

\noindent
\textbf{Role.}

\begin{itemize}
    \item stabilize deep architectures
    \item ensure gradient flow
\end{itemize}

\medskip

\noindent
\textbf{Connection to earlier lectures.}

\begin{itemize}
    \item residuals $\rightarrow$ enable depth (Lecture 20)
    \item normalization $\rightarrow$ stabilize training (Lecture 18)
\end{itemize}

\medskip

\noindent
\textbf{Insight.}  
Transformers integrate multiple stabilization techniques.

\subsection{Computational Characteristics}

\textbf{Advantages.}

\begin{itemize}
    \item full parallelism across sequence
    \item direct modeling of long-range dependencies
\end{itemize}

\medskip

\noindent
\textbf{Cost.}

\begin{itemize}
    \item attention requires $O(T^2)$ computation
\end{itemize}

\medskip

\noindent
\textbf{Tradeoff.}

\begin{itemize}
    \item higher expressivity vs computational cost
\end{itemize}

\medskip

\noindent
\textbf{Insight.}  
Transformers trade efficiency for expressive global interaction.

\subsection{Scaling Laws}

Modern transformer-based models exhibit predictable scaling behavior.

\medskip

\noindent
\textbf{Empirical observation.}

\begin{itemize}
    \item performance improves with:
    \begin{itemize}
        \item model size (parameters)
        \item dataset size
        \item compute
    \end{itemize}
\end{itemize}

\medskip

\noindent
\textbf{Scaling law (informal).}

\[
\text{Loss} \propto N^{-\alpha},
\]

where $N$ represents scale (data, parameters, compute).

\medskip

\noindent
\textbf{Implication.}

\begin{itemize}
    \item larger models trained on more data perform better
\end{itemize}

\medskip

\noindent
\textbf{Insight.}  
Performance follows smooth scaling trends rather than abrupt improvements.

\subsection{Transformers in Practice}

Transformers are the foundation of modern AI systems:

\begin{itemize}
    \item large language models (LLMs)
    \item vision transformers
    \item multimodal models
\end{itemize}

\medskip

\noindent
\textbf{Key property.}

\begin{itemize}
    \item flexible across domains
    \item unified architecture for different data types
\end{itemize}

\medskip

\noindent
\textbf{Insight.}  
Transformers provide a general-purpose architecture for learning representations.

\subsection{Connection to Previous Concepts}

Transformers unify ideas developed throughout the course:

\begin{itemize}
    \item kernel methods $\rightarrow$ attention as learned similarity
    \item deep networks $\rightarrow$ compositional structure
    \item residual connections $\rightarrow$ stable depth
    \item normalization $\rightarrow$ controlled training dynamics
\end{itemize}

\medskip

\noindent
\textbf{Insight.}  
Transformers synthesize multiple strands of machine learning into a single framework.

\subsection{A Unifying Perspective}

Transformers can be viewed as:

\begin{itemize}
    \item \textbf{Interaction networks:} modeling pairwise relationships
    \item \textbf{Adaptive operators:} data-dependent transformations
    \item \textbf{Scalable architectures:} benefiting from scale
\end{itemize}

\medskip

\noindent
\textbf{Core idea.}  
Learning is achieved by iteratively refining representations through global interactions.

\medskip

\noindent
\textbf{Key insight.}  
The success of modern AI arises from combining expressive architectures with large-scale optimization.

\subsection{Outlook}

In this lecture, we studied:

\begin{itemize}
    \item transformer architecture,
    \item positional encoding,
    \item residual and normalization mechanisms,
    \item scaling laws and modern applications.
\end{itemize}

\medskip

\noindent
Transformers represent a major milestone in machine learning.

\medskip

\noindent
In the following part of the course, we connect these ideas to generative modeling, large-scale systems, and theoretical perspectives.

\medskip

\noindent
\textbf{Guiding principle.}  
Modern machine learning systems are shaped by both architectural design and scale.